\chapter{Polynomial potentials and boundary algebras}
\label{ch:polynomial-boundary-algebras}

The equation \(D^2=W\,1\), with \(D\) a matrix and \(W\) a polynomial,
allows an operator whose square is prescribed rather than zero.
Commuting with \(D\), with the appropriate signs, nevertheless gives
a differential on matrices. This elementary construction distinguishes
three questions: which boundary states survive, which polynomial
relations vanish on every boundary, and which relations vanish in the
endomorphisms of a diagonal factorization. Even in one variable the
answers differ.

The value of the potential is part of the question. A critical point
is a point where all first derivatives vanish. The polynomials
\(x^2\) and \(x^2+1\) have the same derivative and the same critical
point. Their boundary algebras behave differently: the first has
nonzero cohomology, whereas every boundary for the second is
contractible. The difference of two values,
\(W(x)-W(y)\), forgets the added constant. A construction made from
that difference must therefore specify how it recovers the chosen
value of the potential.

\section{A differential on boundary operators}
\label{sec:pb-boundary-differential}

Let \(k\) be a field of characteristic zero: every nonzero element has
a multiplicative inverse, and no positive integer multiple of \(1\)
is zero. Let \(S=k[x_1,\ldots,x_n]\) be the ring of finite polynomial
sums in commuting variables. A finite free \(S\)-module is a space of
vectors with polynomial entries and a fixed finite basis. A parity
decomposition \(E=E^{\bar0}\oplus E^{\bar1}\) divides its basis into
even and odd vectors. An even map preserves the two summands; an odd
map exchanges them. The parity \(|A|\) of a homogeneous map is \(0\)
or \(1\), with addition modulo two.

Fix \(W\in S\). A matrix factorization of \(W\) is such a module \(E\)
together with an odd \(S\)-linear map \(D:E\to E\) satisfying
\(D^2=W\,1_E\), where \(1_E\) is the identity map. For two
factorizations \((E,D_E)\) and \((F,D_F)\) of the same \(W\), define
the differential on a homogeneous \(S\)-linear map \(A:E\to F\) by
\begin{equation}
 \delta A=D_FA-(-1)^{|A|}AD_E .
 \label{eq:pb-differential}
\end{equation}
Its square is zero, since expansion gives
\[
 \delta^2A=D_F^2A-AD_E^2=WA-AW=0.
\]
For composable maps \(A:E\to F\) and \(B:F\to G\), direct expansion
also gives
\[
 \delta(BA)=(\delta B)A+(-1)^{|B|}B(\delta A).
\]
A map is closed if its differential is zero and exact if it is the
differential of another map. Closed maps modulo exact maps form the
cohomology of this two-periodic complex. The displayed product rule
shows that composition is defined on cohomology.

These objects, map complexes, and compositions define the
differential graded category \(\operatorname{MF}_{\mathrm{fr}}(S,W)\).
The subscript specifies finite free modules. Its even closed maps
modulo even exact maps are the morphisms in its homotopy category.
All statements below concern this explicitly defined category.

An object is contractible if \(1_E=\delta h\) for an odd endomorphism
\(h\). Then every closed map out of \(E\) is exact: if \(A:E\to F\)
has parity \(p\), the product rule gives
\(\delta((-1)^pAh)=A\). The corresponding assertion for maps into
\(E\) follows from \(\delta(hA)=A\). Thus a contracting homotopy removes
the object and all its cohomology classes of morphisms.

\begin{proposition}[The potential and its derivatives]
\label{prop:pb-potential-primitives}
For every matrix factorization over \(S\),
\[
 \delta(D/2)=W\,1_E,\qquad
 \delta(\partial_iD)=(\partial_iW)\,1_E
 \quad(1\leq i\leq n).
\]
Here \(\partial_i\) denotes entrywise polynomial differentiation with
respect to \(x_i\). Consequently, if polynomials \(a,b_1,\ldots,b_n\)
satisfy \(aW+\sum_i b_i\partial_iW=1\), the odd map
\[
 h=aD/2+\sum_i b_i\partial_iD
\]
contracts \(E\).
\end{proposition}
\begin{proof}
The map \(D\) is odd, so its differential is \(2D^2\).
Differentiating \(D^2=W\,1_E\) gives
\((\partial_iD)D+D(\partial_iD)=(\partial_iW)\,1_E\).
Scalar polynomials commute with all matrices and have differential
zero. Taking the indicated linear combination proves the assertion.
\end{proof}

For \(W=x^2+1\), the identity \(W-(x/2)W'=1\) gives
\[
 h=\tfrac12D-\tfrac{x}{2}\partial_xD,\qquad \delta h=1_E.
\]
This holds for every finite rank. The Jacobi algebra is the quotient
\[
 \operatorname{Jac}(W)=S/(\partial_1W,\ldots,\partial_nW).
\]
For \(W=x^2+1\) it is \(k[x]/(2x)=k\).
It therefore cannot, under these hypotheses,
be the cohomology of boundary operators for the zero value of \(W\).

The ideal \((W,\partial_1W,\ldots,\partial_nW)\) consists of the
polynomial combinations used in the proposition. Its vanishing
requires both \(W=0\) and \(dW=0\). Replacing \(W\) by \(W-c\), for
\(c\in k\), selects the critical points whose value is \(c\).

\section{The complete rank-one calculation}
\label{sec:pb-rank-one}

Over \(S=k[x]\), choose nonzero polynomials \(a,b\) with \(ab=W\).
The rank-one factorization \(K(a,b)\) has one even and one odd basis
vector and matrix
\[
 D=\begin{pmatrix}0&a\\ b&0\end{pmatrix}.
\]
Its endomorphisms comprise all two-by-two polynomial matrices.
The diagonal matrices are even and the off-diagonal matrices are odd.
Formula \eqref{eq:pb-differential} becomes
\begin{align}
 \delta\begin{pmatrix}p&0\\0&q\end{pmatrix}
  &=\begin{pmatrix}0&a(q-p)\\b(p-q)&0\end{pmatrix},
  \label{eq:pb-even-differential}\\
 \delta\begin{pmatrix}0&u\\v&0\end{pmatrix}
  &=(av+bu)\,1.
  \label{eq:pb-odd-differential}
\end{align}

Polynomial division gives, for a nonzero divisor polynomial \(q\),
\(f=cq+r\) with \(\deg r<\deg q\). Repeated division of \(a,b\)
terminates because the nonzero remainder degrees decrease. A
polynomial with leading coefficient one is called monic. The last
nonzero remainder, rescaled to be monic, is their
greatest common divisor \(g\). Reversing the divisions expresses
\(g=\alpha a+\beta b\) for polynomials \(\alpha,\beta\). Thus
the ideal \((a,b)\) is precisely \((g)\). Put
\(a_0=a/g\) and \(b_0=b/g\); these are relatively prime.

\begin{theorem}[Endomorphisms of a rank-one factorization]
\label{thm:pb-rank-one}
Let \(g\) be as above and let
\[
 \eta=\begin{pmatrix}0&a_0\\-b_0&0\end{pmatrix}.
\]
The cohomology algebra of \(\operatorname{End}_S(K(a,b))\) is
\begin{equation}
 H^{\bar0}=S/(g),\qquad H^{\bar1}=(S/(g))[\eta],
 \qquad [\eta]^2=-[a_0b_0].
 \label{eq:pb-rank-one-cohomology}
\end{equation}
The scalar class is represented by a multiple of the identity.
The symbol \([\eta]\) denotes an odd generator commuting with the
even scalars; graded commutativity is not imposed.
The factorization is contractible if and only if \(g=1\).
\end{theorem}
\begin{proof}
Since \(a,b\) are nonzero and \(S\) has no zero divisors, the closed
even matrices in \eqref{eq:pb-even-differential} have \(p=q\).
Their exact scalar coefficients form \((a,b)=(g)\) by
\eqref{eq:pb-odd-differential}.

An odd matrix is closed exactly when \(a_0v+b_0u=0\).
Choose \(\alpha_0,\beta_0\) with
\(\alpha_0a_0+\beta_0b_0=1\). The equation shows that \(a_0\)
divides \(b_0u\); the Bezout identity then shows that \(a_0\)
divides \(u\). Write \(u=a_0h\). Substitution and cancellation of
\(a_0\) give \(v=-b_0h\). Thus every odd cycle is \(h\eta\).
The odd boundaries in \eqref{eq:pb-even-differential} are precisely
the multiples of \(g\eta\). Finally, multiplying the displayed
matrix \(\eta\) by itself gives \(-a_0b_0\,1\).
The identity is exact precisely when \(1\in(g)\), which for monic
\(g\) means \(g=1\).
\end{proof}

For \(a=b=x\), the result is the two-dimensional algebra
\[
 H^\bullet\operatorname{End}(K(x,x))
   =k\,1\oplus k[\eta],\qquad [\eta]^2=-1.
\]
This is the Clifford algebra on one odd generator with the stated
square. Over \(\mathbb C\), multiplication of the generator by \(i\)
changes its square to \(1\). The differential graded endomorphism
algebra itself still contains all polynomial matrices. Its
two-dimensional cohomology records even and odd endomorphism
classes of one object; it is not a count of distinct objects.

More generally, let \(W=x^m\) and \(1\leq j\leq m-1\).
For \(K(x^j,x^{m-j})\), put
\(e=\min(j,m-j)\). Then
\[
 H^\bullet\operatorname{End}(K(x^j,x^{m-j}))
   =\bigl(k[x]/(x^e)\bigr)\,1
    \oplus\bigl(k[x]/(x^e)\bigr)[\eta],
 \qquad [\eta]^2=-[x^{m-2e}].
\]
The common power removed from the two differentials determines the
remaining scalar relations, while the excess power determines the
odd product.

\section{Reduction of every finite rank}
\label{sec:pb-every-rank}

The rank-one calculation determines all finite free objects in one
variable because polynomial division also diagonalizes matrices.
An invertible polynomial matrix means a matrix whose inverse again
has polynomial entries. Interchanging two rows, adding a polynomial
multiple of one row to another, and multiplying a row by a nonzero
constant are invertible operations; the same holds for columns.

\begin{lemma}[Polynomial diagonal reduction]
\label{lem:pb-smith}
If a square matrix \(A\) over \(k[x]\) has nonzero determinant, there
are invertible polynomial matrices \(U,V\) such that
\[
 UAV=\operatorname{diag}(a_1,\ldots,a_r),\qquad
 a_1\mid a_2\mid\cdots\mid a_r,
\]
with each \(a_i\) monic and nonzero.
\end{lemma}
\begin{proof}
Move a nonzero entry into the upper left corner. Divide each entry
below it by the corner entry, subtract the quotient times the first
row, and, if the remainder is nonzero, move that remainder into the
corner. Its degree is smaller. Perform the corresponding column
operations on the first row. Each nonzero remainder reduces the
corner degree, so after finitely many such reductions the corner
divides every entry in its row and column. Clear that row and column.

If the corner fails to divide an entry of the remaining block, add
the row containing that entry to the first row. The corner is
unchanged, and polynomial division of the new first-row entry
produces a nonzero remainder of smaller degree. Repeat the reduction.
This can occur only finitely often. The resulting corner divides
every entry of the remaining block. Apply the same argument to that
block. Its subsequent entries remain multiples of the first corner
because all row and column operations have polynomial coefficients.
Induction gives the divisibility chain. Nonzero determinant excludes
zero diagonal entries. Rescaling rows makes them monic.
\end{proof}

\begin{theorem}[Finite free factorizations in one variable]
\label{thm:pb-all-ranks}
Let \(0\ne W\in k[x]\). Every object of
\(\operatorname{MF}_{\mathrm{fr}}(k[x],W)\) is isomorphic, before
taking cohomology, to a finite direct sum of \(K(a,W/a)\), where
\(a\) is a monic divisor of \(W\).
\end{theorem}
\begin{proof}
Write \(D=\left(\begin{smallmatrix}0&A\\B&0\end{smallmatrix}\right)\).
The equations \(AB=W\,1\) and \(BA=W\,1\) imply, over the field
of rational functions \(k(x)\), that \(A,B\) are inverse up to the
nonzero scalar \(W\). Their source and target have equal dimension,
and \(A\) has nonzero determinant.

Apply Lemma~\ref{lem:pb-smith}. An even change of basis conjugates
\(D\) to the matrix with off-diagonal entries
\(A'=UAV=\operatorname{diag}(a_i)\) and
\(B'=V^{-1}BU^{-1}\). The equation \(A'B'=W\,1\) forces every
off-diagonal entry of \(B'\) to be zero, because each \(a_i\)
is nonzero. Its diagonal entries are \(W/a_i\). These are
polynomials since \(B'\) is a polynomial matrix. The resulting
factorization is the stated direct sum.
\end{proof}

This reduction also gives a finite description of all morphism
complexes. For divisors \(a,c\) of \(W\), put \(b=W/a\), \(d=W/c\).
A map from \(K(a,b)\) to \(K(c,d)\) has differential
\[
 \delta\begin{pmatrix}p&0\\0&q\end{pmatrix}
   =\begin{pmatrix}0&cq-pa\\dp-qb&0\end{pmatrix},
 \qquad
 \delta\begin{pmatrix}0&u\\v&0\end{pmatrix}
   =\begin{pmatrix}cv+ub&0\\0&du+va\end{pmatrix}.
\]
Composition is ordinary matrix multiplication. There are finitely
many monic divisors of \(W\). Their blocks, these map complexes, and
finite direct sums therefore describe the entire category defined
in Section~\ref{sec:pb-boundary-differential}, without replacing
its differential by its cohomology.

To compare scalar relations across all objects, define
\[
 I_{\mathrm{obj}}(W)
 =\{s\in k[x]:s\,1_E\text{ is exact for every finite free
 factorization }E\text{ of }W\}.
\]
It is an ideal: sums of primitives and multiplication of a primitive
by a polynomial give primitives for the corresponding sums and
multiples.

A nonconstant polynomial is irreducible if it is not a product of
two nonconstant polynomials. These polynomials determine the common
scalar ideal one multiplicity at a time.

\begin{theorem}[Scalars vanishing on every object]
\label{thm:pb-object-ideal}
Factor \(W=c\prod_{\nu=1}^q p_\nu^{m_\nu}\), where \(c\in k^\times\),
the \(p_\nu\) are distinct monic irreducible polynomials, and
\(m_\nu\geq1\). Then
\begin{equation}
 I_{\mathrm{obj}}(W)
   =\left(\prod_{\nu=1}^q
                 p_\nu^{\lfloor m_\nu/2\rfloor}\right).
 \label{eq:pb-object-ideal}
\end{equation}
For a nonzero constant \(W\), the empty product is \(1\).
\end{theorem}
\begin{proof}
Repeated polynomial division proves existence of factorization into
irreducibles by induction on degree. Bezout's identity shows that
an irreducible dividing a product divides one factor; this proves
uniqueness of the multiplicities used here.

For \(a=\prod_\nu p_\nu^{j_\nu}\), its complementary factor \(W/a\)
has exponents \(m_\nu-j_\nu\). Theorem~\ref{thm:pb-rank-one} gives
the scalar ideal
\((\prod_\nu p_\nu^{\min(j_\nu,m_\nu-j_\nu)})\) on this block.
A scalar contracts on every direct sum if it contracts on every
block. Conversely, projecting a contracting homotopy on a direct
sum to one block gives a contracting homotopy there. Thus
Theorem~\ref{thm:pb-all-ranks} identifies \(I_{\mathrm{obj}}\) with
the intersection of these principal ideals over all divisors \(a\).

Membership in that intersection requires the exponent of each
\(p_\nu\) to be at least
\[
 \max_{0\leq j\leq m_\nu}\min(j,m_\nu-j)
 =\lfloor m_\nu/2\rfloor .
\]
Unique factorization gives
\eqref{eq:pb-object-ideal}.
\end{proof}

For \(W=x^m\), the common scalar ideal is
\((x^{\lfloor m/2\rfloor})\). A primitive may depend on the
object. The definition does not require those choices to commute
with maps between objects. The next calculation measures scalar
relations on a different, explicitly specified algebra.

\section{The diagonal factorization}
\label{sec:pb-diagonal}

Let \(W\in k[x]\) be nonconstant, introduce a second independent
variable \(y\), and write
\[
 T=k[x,y],\qquad r=x-y,\qquad
 g=\frac{W(x)-W(y)}{x-y}.
\]
The quotient is a polynomial: for a monomial,
\((x^j-y^j)/(x-y)=\sum_{\ell=0}^{j-1}x^{j-1-\ell}y^\ell\).
The diagonal factorization is the rank-one object over \(T\)
\[
 D_\Delta=\begin{pmatrix}0&r\\g&0\end{pmatrix},\qquad
 D_\Delta^2=(W(x)-W(y))\,1.
\]
Its differential graded endomorphism algebra will be denoted
\(B_W=(\operatorname{End}_T(T^{\bar0}\oplus T^{\bar1}),\delta)\).
The name refers to the equation \(r=0\), which identifies the two
coordinates. The definition of \(B_W\) does not identify it with
the endomorphisms of an identity functor on another category.

\begin{theorem}[The diagonal endomorphism algebra in one variable]
\label{thm:pb-diagonal}
Scalar matrices induce an isomorphism of algebras
\[
 H^{\bar0}(B_W)=T/(r,g)\cong k[x]/(W'(x)),\qquad
 H^{\bar1}(B_W)=0 .
\]
The kernel of the left scalar map \(k[x]\to H^{\bar0}(B_W)\)
is \((W')\).
\end{theorem}
\begin{proof}
The even cycles are scalar matrices, and the even boundaries
have coefficients in \((r,g)\), by
\eqref{eq:pb-even-differential}--\eqref{eq:pb-odd-differential}.
Modulo \(r\), the monomial formula for \(g\) gives \(g=W'(x)\).
Characteristic zero and nonconstancy imply \(W'\ne0\).

An odd cycle with entries \(u,v\) satisfies \(rv+gu=0\).
Reducing modulo \(r\) gives \(W'(x)u(x,x)=0\) in the integral
domain \(k[x]\). Hence \(u=rh\) for some \(h\in T\);
division by the monic polynomial \(x-y\) justifies this last
assertion. Substitution gives \(v=-gh\). This odd cycle is
\(\delta\operatorname{diag}(0,h)\), so odd cohomology vanishes.
The scalar representatives multiply as their coefficients do.
\end{proof}

For \(W=x^m\), \(m\geq2\), this diagonal has left scalar kernel
\((x^{m-1})\), while all finite free boundary objects have common
scalar kernel \((x^{\lfloor m/2\rfloor})\). At \(m=3\), for example,
\(x\) acts by an exact map on every boundary object but is nonzero
in \(H^{\bar0}(B_W)=k[x]/(x^2)\). The two computations compare
specified algebras and ideals. To interpret the diagonal algebra
as categorical Hochschild cochains one must further construct its
action as an identity kernel, with the chosen potential value retained.

The derivative homotopy is already explicit. Put
\[
 H=\partial_xD_\Delta
   =\begin{pmatrix}0&1\\\partial_xg&0\end{pmatrix}.
\]
Proposition~\ref{prop:pb-potential-primitives}, applied to
\(W(x)-W(y)\) with \(y\) held fixed, gives
\[
 \delta H=W'(x)\,1,\qquad
 H^2=(\partial_xg)\,1,\qquad
 (\partial_xg)(x,x)=\tfrac12W''(x).
\]
The last identity follows by differentiating the monomial sum for
\(g\): its coefficients add to \(j(j-1)/2\) for \(W=x^j\).
Thus the Hessian term arises from multiplying an actual primitive.
The same primitive is an odd cycle only when \(W'=0\) in the
chosen scalar ring.

Right differentiation gives a different sign. The odd map
\(-\partial_yD_\Delta\) has differential \(W'(y)\,1\);
its square on the diagonal is \(-W''(y)/2\).
Which scalar action is retained determines this sign.

\section{Separating the critical values}
\label{sec:pb-critical-values}

Adding a constant to \(W\) leaves \(B_W\) unchanged. Theorem
\ref{thm:pb-diagonal} therefore detects all critical points of \(W\),
including those away from its zero fibre. In one variable they can
be separated inside the Jacobi algebra by explicit polynomial
idempotents.

Assume for this section that every nonconstant polynomial over \(k\)
splits into linear factors; such a field is called algebraically
closed. Assume also that \(\deg W\geq2\). Write
\[
 W'(x)=a\prod_{\alpha\in A}(x-\alpha)^{e_\alpha},
 \qquad a\in k^\times,\quad e_\alpha\geq1,
\]
with \(A\) the finite set of distinct critical points.
The set of critical values is
\(V=\{W(\alpha):\alpha\in A\}\).
Put \(J=k[x]/(W')\), and denote the class of a polynomial \(f\)
by \([f]\).

\begin{lemma}[The potential in the Jacobi algebra]
\label{lem:pb-critical-projectors}
There is an algebra isomorphism
\[
 J\longrightarrow
 \prod_{\alpha\in A} k[x]/((x-\alpha)^{e_\alpha})
\]
given by reduction of a polynomial in each factor.
In the factor indexed by \(\alpha\), the class of \(W\) is
the constant \(W(\alpha)\).
\end{lemma}
\begin{proof}
Powers of two distinct linear polynomials are relatively prime.
Bezout's identity therefore gives elements congruent to \(1\)
modulo one such power and to \(0\) modulo all the others.
Multiplying each desired component by the corresponding element
proves surjectivity of the reduction map. Its kernel is the
intersection of the power ideals, which is their product by
unique factorization. This product generates \((W')\), proving
the isomorphism.

Expand \(W(\alpha+z)=\sum_j c_jz^j\).
The multiplicity \(e_\alpha\) of the zero of \(W'\), and
characteristic zero, imply \(c_1=\cdots=c_{e_\alpha}=0\).
Thus \(W-W(\alpha)\) is divisible by
\((x-\alpha)^{e_\alpha+1}\), and has zero class in the indicated
factor.
\end{proof}

For \(c\in V\) define
\[
 P_c(X)=\prod_{\substack{d\in V\\d\ne c}}\frac{X-d}{c-d},
 \qquad p_c=P_c(W(x)),\qquad e_c=[p_c]\in J.
\]
Every denominator is a nonzero element of \(k\).
An empty product is \(1\). The lemma gives
\begin{equation}
 e_ce_d=\begin{cases}e_c&c=d,\\0&c\ne d,\end{cases}
 \qquad \sum_{c\in V}e_c=1,\qquad
 [W]e_c=c\,e_c.
 \label{eq:pb-critical-projectors}
\end{equation}
Indeed, on the factor for \(\alpha\), \(e_c\) has value \(1\)
if \(W(\alpha)=c\), and \(0\) otherwise. Each factor belongs to
exactly one critical value.

The zero-value algebra \(J/([W])\) is \(e_0J\) if \(0\in V\),
and is zero otherwise. On every component with \(c\ne0\),
multiplication by \([W]\) is multiplication by the invertible
scalar \(c\); on the zero component it is zero.
This is an algebraic separation of the critical values.

The projectors also admit explicit homotopies before taking
cohomology. By \eqref{eq:pb-critical-projectors} the polynomial
\(p_cp_d-\delta_{cd}p_c\) is divisible by
\(W'\), where \(\delta_{cd}\) is \(1\) when \(c=d\) and \(0\)
otherwise. Define its polynomial quotient by
\[
 q_{cd}=\frac{p_cp_d-\delta_{cd}p_c}{W'}.
\]
Using the matrix \(H\) above gives
\[
 \delta(q_{cd}H)=(p_cp_d-\delta_{cd}p_c)\,1.
\]
Completeness holds exactly: \(\sum_cP_c(X)-1\) has degree less
than \(|V|\) and vanishes at all \(|V|\) distinct elements of \(V\),
so it is zero by repeated polynomial division. Hence
\(\sum_cp_c=1\) already in \(k[x]\). Orthogonality and idempotence
hold by the displayed homotopies in \(B_W\).

\begin{example}[Two critical values]
\label{ex:pb-two-values}
Take \(W=x^3-3x+2=(x-1)^2(x+2)\).
Then \(W'=3(x^2-1)\), the critical points are \(1,-1\), and
their values are \(0,4\). In \(J=k[x]/(x^2-1)\),
\[
 [W]=2-2[x],\qquad
 e_0=1-[W]/4=(1+[x])/2,\qquad
 e_4=[W]/4=(1-[x])/2.
\]
Both are nonzero and their sum is one.
The zero-value quotient is \(J/([W])=k\), whereas
the whole diagonal cohomology has dimension two.
Theorem~\ref{thm:pb-object-ideal} gives
\(I_{\mathrm{obj}}(W)=(x-1)\). In particular, the critical point
\(-1\) contributes to the diagonal Jacobi algebra but supplies
no additional scalar component on the boundary objects for \(W=0\).
\end{example}

\section{Parity and the remaining comparison}
\label{sec:pb-parity-comparison}

A parity shift exchanges \(E^{\bar0}\) and \(E^{\bar1}\) and changes
the sign of \(D\); this operation is denoted \([1]\). Applying it
twice restores the original parity and the original differential,
so \([2]\) is the identity in the two-periodic category.
An integer cohomological grading, by contrast, consists of separate
spaces in every integer degree and a differential raising degree
by one. Such a grading contains data absent from a parity
decomposition. In particular, specifying a parity shift does not
distinguish an integer \(n\) from \(n-2\).

The boundary calculation and the diagonal calculation can be
placed in one diagram of scalar maps:
\[
\begin{array}{ccc}
 k[x] &\longrightarrow& H^{\bar0}(B_W)=k[x]/(W')\\
 \big\downarrow && \big\downarrow\\
 k[x]/I_{\mathrm{obj}}(W)&& k[x]/(W,W').
\end{array}
\]
The vertical arrows and the top arrow are the displayed quotient
maps. No bottom arrow is asserted. For \(W=x^3\), such an arrow
preserving the specified scalar actions would send the zero class
of \(x\) in \(k[x]/(x)\) to the nonzero class of \(x\) in
\(k[x]/(x^2)\), which is impossible.

The algebra of one boundary, the relations vanishing on every
boundary, and the endomorphisms of the diagonal are therefore
distinct inputs to a comparison. A categorical identification
requires a diagonal action on the boundary category and a proof
that this action represents its identity functor. A chiral
realization additionally requires local operations depending on
positions, their differential, and the identities governing
collisions. The finite polynomial matrices specify the boundary
input and its homotopies; they do not supply those additional
position-dependent operations.
